\section{Related Work}

Long‑running autonomous systems have been studied in several communities.
Autonomic computing~\cite{kephart2003autonomic} introduced the notion of
self‑configuring, self‑healing, self‑optimising, and self‑protecting
systems, but it remained largely at the design‑pattern level.  The EU
FP7 project ASCENS~\cite{wirsing2012ascens} developed formal methods for
self‑aware and self‑expressive ensembles, yet its focus was on scientific
illustrations rather than a reusable, lightweight core.

Event‑driven architectures are the backbone of many distributed systems
(e.g., Apache Kafka, AWS EventBridge).  However, these systems are
engineered for high throughput and durability in commercial applications,
not for introspective, ethically governed knowledge systems.  The \Core{}
event bus borrows patterns such as persistent logging and dead‑letter
queues but adds first‑class support for event sourcing and replay that
is essential for auditability.

Adaptive thresholds draw on control theory (PID controllers) and
meta-learning~\cite{finn2017maml}.  Our module couples a Kalman filter
with a \texttt{ComplexityProvider} interface, making it adaptable to
any system that reports its complexity and unresolved gaps.

Chaos detection in running systems is a central theme of chaos
engineering~\cite{basiri2016chaos}.  Netflix's Simian Army and similar
tools inject failures to test resilience.  The \Core{} chaos detector is
inspired by these ideas but operates passively: it monitors internal
metrics (error epsilon, divergence rate, Lyapunov exponent) and issues
graded interventions rather than injecting faults.

Finally, the thermodynamic cost model is rooted in Landauer's
principle~\cite{landauer1961} and the growing literature on Green AI.
Recent work has proposed carbon‑aware job scheduling for data
centres~\cite{acun2023carbon}; we adapt this idea to the granularity of
individual computational structures within a single agent.

\section{Architecture Overview}

The \Core{} follows a strictly event-driven, modular design (Figure~\ref{fig:core_arch}).  
All communication between the four main subsystems---EventBus, AdaptiveThresholds,
ChaosDetector, and ThermodynamicCost---passes through the central event bus.  
No component holds a direct reference to another; they interact only via typed
messages published on named channels.  This discipline enables (1) transparent
logging of every decision, (2) hot-swapping of individual modules, and
(3) deterministic replay of system behaviour for debugging or forensics.

\begin{figure}[htbp]
\centering
\begin{tikzpicture}[
    node distance=1.8cm and 2.0cm,
    module/.style={draw, rounded corners, minimum width=2.8cm, minimum height=1.0cm, align=center, font=\small},
    bus/.style={draw, fill=gray!10, minimum width=8cm, minimum height=1.0cm, align=center, font=\small\bfseries},
    arrow/.style={-{Stealth[length=2mm]}, thick}
]
\node[bus] (bus) {EventBus (asyncio / Redis)};
\node[module, fill=blue!10, above=of bus, xshift=-3cm] (adapt) {Adaptive\\Thresholds};
\node[module, fill=red!10, above=of bus, xshift=0cm] (chaos) {Chaos\\Detector};
\node[module, fill=green!10, above=of bus, xshift=3cm] (thermo) {Thermodynamic\\Cost};
\node[module, fill=orange!10, below=of bus] (logger) {Persistent\\Event Store};

\draw[arrow] (adapt) -- (bus);
\draw[arrow] (chaos) -- (bus);
\draw[arrow] (thermo) -- (bus);
\draw[arrow] (bus) -- (logger);
\end{tikzpicture}
\caption{Architecture of the \Core{}.  All modules communicate exclusively through the event bus; events are durably logged.}
\label{fig:core_arch}
\end{figure}

The \Core{} is lightweight: the entire implementation is contained in five
Python files (event\_bus.py, adaptive\_thresholds.py, chaos\_detection.py,
thermodynamic\_cost.py, and the shared config.py) totaling approximately 
2,800 lines of code.  It has no mandatory external dependencies beyond the
Python standard library, psutil, and numpy; optional integrations (Redis,
Qiskit, pynvml) are activated only when the corresponding hardware or service
is available.

The remainder of this section briefly describes each module; detailed
technical exposition follows in Sections~3--6.